\documentclass[10pt]{article}
\usepackage[margin=20mm]{geometry}
\usepackage{amsmath,amssymb,mathtools,fontspec,unicode-math,xurl,hyperref}
\IfFontExistsTF{DejaVu Serif}{\setmainfont{DejaVu Serif}}{\setmainfont{Latin Modern Roman}}
\setmathfont{latinmodern-math.otf}
\setlength{\parindent}{0pt}\setlength{\parskip}{6pt}\sloppy
\title{Initial scientific review: M25-0042
Hermitian K-theory for stable \(\infty\)-categories II}
\author{}
\date{25 October 2026}
\begin{document}
\maketitle
Abstract: No substantive issue was established in the declared, claim-focused coverage. The review checks the distinction between additive and localising invariants, the realisation step in additivity, the Bott–Genauer shifts, the passage from negative to all homotopy groups, and genuine equivariant periodicity. A suspension-placement inconsistency in the final display on printed page 338 is separately documented, with its derivation and limitations. This is not certification of every argument in the 252-page paper.\par
\section*{Identity, version and disposition}
The reviewed object is the supplied publisher PDF of Baptiste Calmès, Emanuele Dotto, Yonatan Harpaz, Fabian Hebestreit, Markus Land, Kristian Moi, Denis Nardin, Thomas Nikolaus and Wolfgang Steimle, Hermitian K-theory for stable \(\infty\)-categories II: Cobordism categories and additivity, Acta Mathematica 235 (2025), 149–400. All locations below use printed journal pages, with one-based PDF pages in parentheses. Printed page equals PDF page plus 148.\par
The title page prints DOI `10.4310/ACTA.2025.v235.n2.a1`. The frozen assignment instead records `10.4310/acta.2025.n235.n2.a1`. The volume-prefix discrepancy is preserved as metadata, not silently rewritten or counted as a scientific defect. The source is `papers/0042/0042.pdf`; its SHA-256 is\par
> 949d7a12e4a7546c10acec150efdc69bb
> 2f95244d965165c1e3ca1a5c51bca7d.\par
The complete supplied manifest was parsed and every bound file checked during intake. That operation is not scientific review. The assigned identity was not replaced by an arXiv revision. A current arXiv metadata check identifies v5 dated 11 April 2025; no assertion of byte or textual equivalence with this publisher PDF is made.\par
Disposition: No Issue Found within the substantive coverage below. There are zero confirmed Major and zero confirmed Minor findings. Observation X01 is excluded from those counts as an isolated display inconsistency: the relevant preceding corollary supplies the correctly shifted conclusion. This classification does not assert author intent or independent human acceptance.\par
\section*{Claims, hypotheses and independent checks}
\subsection*{Additivity is a condition on squares}
Definition 1.6.3, p. 203 (PDF 55), requires a reduced functor, a target admitting finite limits, and preservation of the appropriate class of squares. Proposition 1.6.4, p. 204 (PDF 56), permits replacement by sequences when the target is stable. The text explicitly retains the stronger square condition for a space-valued functor.\par
This distinction is necessary in the proof of Theorem 2.5.3 on p. 259 (PDF 111): one square recognises cocartesian arrows and another supplies their lifts. Merely assuming that a space-valued functor preserves sequences would not justify applying it to these squares. The actual definition supplies the stronger assumption, so this potential objection is rejected. The split condition on the projection is also explicit in Theorem 2.5.3, p. 249 (PDF 101), and the failure of an unrestricted nonsplit extension is acknowledged in Remark 2.5.4(ii), p. 250 (PDF 102).\par
The terminal/zero case is consistent: reducedness forces the invariant of the zero Poincaré category to be terminal; for a stable target this is the zero object. Thus no extra summand occurs when one term of a split sequence is zero.\par
\subsection*{Why realisation preserves the square at issue}
The proof of Theorem 2.5.1 on pp. 250–251 (PDF 102–103) does not use an unrestricted assertion that geometric realisation preserves pullbacks. It uses a bicartesian fibration and the straightening functor into categories with adjoints.\par
Here is the relevant independent implication. If \(L\dashv R\), its unit and counit induce homotopies \[|\mathrm{id}|\simeq |R||L|,\qquad |L||R|\simeq|\mathrm{id}|.\] Hence \(|L|\) is an equivalence. The straightening transition functors therefore become equivalences after groupoid completion and factor through the groupoid completion of the base. Unstraightening this space-valued local system commutes with base change; this is exactly the restricted pullback needed on p. 250.\par
The last step of Theorem 2.5.3, p. 259, also uses the self-opposite symmetry of the hermitian \(Q\)-construction to obtain the cartesian half from the cocartesian half. The recognition step is compatible with the stronger square-based definition of additivity. The full horn and Poincaré-structure calculations in Lemmas 2.6.3–2.6.7 were not all reconstructed; the preceding implication is checked conditional on those lemmas, rather than presented as a new proof of the whole additivity theorem.\par
\subsection*{Bott–Genauer rotation and the \(U,V\) conventions}
Corollary 4.3.1, p. 317 (PDF 169), identifies \(\mathrm{GW}(\mathrm{Met}(\mathcal C,\mathord{\mathrm{Q}}))\) with \(\mathrm{GW}(\mathrm{Hyp}(\mathcal C))\) and then with \(K(\mathcal C)\). The resulting sequence on p. 318 (PDF 170) is \[\mathrm{GW}(\mathcal C,\mathord{\mathrm{Q}}^{[-1]})\longrightarrow K(\mathcal C)
 \longrightarrow\mathrm{GW}(\mathcal C,\mathord{\mathrm{Q}}).\] Write \(G_j=\mathrm{GW}(\mathcal C,\mathord{\mathrm{Q}}^{[j]})\). With the paper’s definitions, \[U(\mathord{\mathrm{Q}}^{[2]})=\operatorname{fib}(K\to G_2)=G_1.\] Applying the sequence to \(\mathord{\mathrm{Q}}^{[1]}\) gives \(G_0\to K\to G_1\); rotation in spectra gives \[V(\mathord{\mathrm{Q}})=\operatorname{fib}(G_0\to K)=\Omega G_1.\] Consequently \(U(\mathord{\mathrm{Q}}^{[2]})\simeq\Sigma V(\mathord{\mathrm{Q}})\), exactly Corollary 4.3.3. Neither a shift by \(-2\) nor an unshifted \(U\simeq V\) follows. This checks the stated rotation; the equivalence between metabolic and hyperbolic theories remains dependent on the cited constructions.\par
\subsection*{Bordification and the fundamental fibre sequence}
Theorem 4.4.11, p. 329 (PDF 181), compares \(\mathrm{GW}^{\mathrm{bord}}\) with \(\mathrm{L}\). Its reduction from negative to all degrees is valid provided the natural negative-degree comparison holds for every shifted Poincaré structure, as asserted through Corollaries 3.6.17 and 4.4.9. For any integer \(i\), choose \(k>i\). For a bordism-invariant functor \(F\), \[F(\mathcal C,\mathord{\mathrm{Q}})\simeq \Sigma^k F(\mathcal C,\mathord{\mathrm{Q}}^{[-k]}),\qquad
 \pi_iF(\mathcal C,\mathord{\mathrm{Q}})=\pi_{i-k}F(\mathcal C,\mathord{\mathrm{Q}}^{[-k]}).\] The right-hand degree is negative. Applying this to both functors and to their natural comparison proves the claimed extension of the isomorphism. Negative-degree agreement for just one unshifted category would be insufficient; the quantification over shifted inputs is essential.\par
Corollary 4.4.13, p. 330 (PDF 182), identifies a cartesian square with bottom row \(K^{hC_2}\to K^{tC_2}\). Taking horizontal fibres yields \(K_{hC_2}\to\mathrm{GW}\to\mathrm{L}\), because the fibre of the norm cofibre map is the homotopy-orbit spectrum. The orbit and fixed-point terms must not be interchanged.\par
For Corollary 4.4.14, the displayed composite from \(K[1/2]_{hC_2}\) through \(\mathrm{GW}[1/2]\) to \((K[1/2])^{hC_2}\) is the norm. Its equivalence after inverting 2 gives a retraction and hence the splitting. The proof uses a map to the fixed points of the localised spectrum; it does not require an unjustified blanket commutation of localisation with arbitrary limits.\par
\subsection*{Localisation, idempotent completion, and a rejected extraction alarm}
Corollary 4.5.1 on p. 330 is Verdier localisation. Examples 1.6.9(iii)–(iv), p. 206 (PDF 58), distinguish connective \(K\), nonconnective \(K\), and connective \(K\) after idempotent completion. The introduction’s claim therefore does not assert that connective \(K\) sends every idempotent-completed quotient sequence to a fibre sequence.\par
Remark 4.5.2, p. 331 (PDF 183), retains involution-closed subgroups \(c\subseteq d\subseteq K_0(R)\) and explicitly cites paper IV for the cofinality squares. Direct page inspection establishes that its Grothendieck–Witt remainder is \[H(d/c)^{hC_2},\] with a superscript. The text extraction obscured this distinction. For an abelian group \(A\) with \(C_2\)-action, \(H(A)^{hC_2}\) has \(\pi_{-j}=H^j(C_2;A)\) for \(j\geq0\) and no positive homotopy groups. The long exact sequence consequently gives isomorphisms in positive degrees and injectivity in degree zero, as stated. An objection based on positive group homology would instead concern \(H(A)_{hC_2}\) and is not an objection to the printed paper. The cited cofinality theorem in paper IV was not independently proved here.\par
\subsection*{Genuine periodicity and its range}
Theorem 3.7.5, p. 311 (PDF 163), identifies the rotated metabolic sequence with the tensor of a genuine spectrum by \[\mathbb S[C_2]\longrightarrow\mathbb S\longrightarrow S^\sigma.\] Thus \(F^{\mathrm{ghyp}}(\mathord{\mathrm{Q}}^{[-1]})\simeq
S^{\sigma-1}F^{\mathrm{ghyp}}(\mathord{\mathrm{Q}})\), and replacing \(\mathord{\mathrm{Q}}\) by \(\mathord{\mathrm{Q}}^{[1]}\) gives the \(1-\sigma\) shift in Corollary 4.6.3, p. 336 (PDF 188). The representation \(1-\sigma\) has underlying dimension zero and fixed dimension one. These two tests give unchanged nonequivariant suspension degree for \(K\), and suspension degree one for geometric fixed points \(\mathrm{L}\). They do not imply ordinary one-periodicity of connective \(K\).\par
Corollary 4.6.4, p. 337 (PDF 189), retains the \(n\sigma\)-orientation, the involution-closed subgroup, and an additional connectivity hypothesis for the truncated structures \(\mathord{\mathrm{Q}}^{\geq m}\). For \(n=2\) it gives \[\mathrm{KR}(\mathord{\mathrm{Q}}_{-M}^{\geq m+1})\simeq
 S^{2-2\sigma}\mathrm{KR}(\mathord{\mathrm{Q}}_M^{\geq m}).\] Taking \(m=0,1\) gives the genuine symmetric/even/quadratic chain with the correct sign change of \(M\). The homotopical structures \(\mathord{\mathrm{Q}}^s,\mathord{\mathrm{Q}}^q\) have the stated period; these should not be confused with an assertion that every genuine truncated structure is four-periodic without changing its form parameter.\par
\section*{Excluded observation X01: final display on page 338}
Location. The last display of §4.6, p. 338 (PDF 190), after Corollary 4.6.7, prints \[\Sigma^4 L_{\mathrm{gq}}\simeq
 \Sigma^2 L_{\mathrm{ge},-M}\simeq L_{\mathrm{gs}},
 \tag{X01-printed}\] where the suppressed category is \(\mathcal D^c(R)\) and the other two coefficients are \(M\).\par
Exact algebraic check. Corollary 4.6.5 on p. 337 instead gives \[S^{2\sigma-2}\mathrm{KR}_{\mathrm{gq}}\simeq
 \mathrm{KR}_{\mathrm{ge},-M}\simeq S^{2-2\sigma}\mathrm{KR}_{\mathrm{gs}}.\] Geometric fixed points send \(a+b\sigma\) to degree \(a\). Therefore \[\Sigma^{-2}L_{\mathrm{gq}}\simeq
 L_{\mathrm{ge},-M}\simeq\Sigma^2L_{\mathrm{gs}},\] and suspension by two gives \[L_{\mathrm{gq}}\simeq
 \Sigma^2L_{\mathrm{ge},-M}\simeq\Sigma^4L_{\mathrm{gs}}.
 \tag{X01-derived}\] This agrees with the stabilisation expression on p. 156 (PDF 8). In particular, the printed final chain is not the chain obtained by the stated derivation. Reversing the direction of an equivalence does not keep a suspension on the same side: from \(A\simeq\Sigma^2B\) one gets \(B\simeq\Sigma^{-2}A\).\par
Counterevidence and impact. The immediately preceding genuine-spectrum corollary, the general shift theorem, and the introduction all supply the consistent direction. The final display is not used in the inspected proof of those results. Remark 4.6.8(i), printed p. 339 (PDF 191), explicitly states that genuine \(L\)-spectra are not periodic in general and refers to a computation of \(L(\mathcal D^p(\mathbb Z),\mathord{\mathrm{Q}}^{\mathrm{gs}})\) in paper III. The printed middle equivalence \(\Sigma^2L_{\mathrm{ge},-M}\simeq L_{\mathrm{gs}}\), combined with \(L_{\mathrm{ge},-M}\simeq\Sigma^2L_{\mathrm{gs}}\) from Corollary 4.6.5, would force \(\Sigma^4L_{\mathrm{gs}}\simeq L_{\mathrm{gs}}\). The adjacent nonperiodicity statement supplies concrete counterevidence to interpreting the reversed chain as an additional periodicity theorem. The cited computation over \(\mathbb Z\) was not independently reconstructed here; no new ring-level calculation or refutation of the main theorem is claimed. The evidence supports correcting the placement of suspensions in this summary display. Because only this local display inconsistency has been established, it is recorded outside the substantive Major/Minor counts under the task’s exclusion of typographical errors alone. No claim about the author’s intended keystroke is needed.\par
\section*{Coverage, evidence and limits}
The text actually read comprises PDF pages 1–2, 7–14, 54–58, 100–103, 110–111, 160–163, 169–170, 181–183 and 188–190: 33 of 252 pages. Reading includes introductory and contextual pages, not 33 pages of fully reconstructed proofs. Actual source-page images were inspected for PDF pages 1, 7, 55, 56, 58, 101, 102, 111, 163, 170, 181, 182, 183, 188, 189 and 190. They bind the identity, hypotheses, decisive formulas and the rejected extraction alarm. Focused acceptance additionally inspected Remark 4.6.8(i), printed p. 339 (PDF 191), for the counterevidence concerning the final display. This additional targeted reading does not expand the initial-review page count, and the explicit computation in paper III was not independently checked.\par
The full surgery construction, all horn fillings, all universal-property arguments, the classification and recollement appendices, comparison to earlier models, and the full ring and parametrised-spectrum applications were not independently reconstructed. Dependencies in paper I, paper IV, and the cited foundational literature were not exhaustively audited. These omissions limit the review; they are not findings against the authors. No computational source repository or separate supplement was acquired in this manifest, and no author code was executed. No numerical experiment is relevant to the categorical checks above; the displayed calculations are independent formal derivations.\par
Current primary-source metadata was inspected at https://arxiv.org/abs/2009.07224 and https://www.math.uni-bielefeld.de/\textasciitilde{}hebestreit/. The latter lists this paper and separately labels an erratum for another paper, On the homotopy type of L-spectra of the integers. That erratum was not attributed to M25-0042. A fetch of the Calmès publication page timed out. No comprehensive absence-of-errata claim is made.\par
The evidence directory preserves the extracted source, page images, I/O log, version-check record and the exact compilation log. The initial native-cloud build is preserved. Focused acceptance compiles the current exact source with the existing authorized Tectonic runtime and records its dependency hashes and actual page inspections before setting delivery readiness. Human scientific adjudication remains pending.\par
\end{document}
